\section{A uniform quantitative amplification theorem}
\label{amp:sec}

The amplification argument of \cite{OpenAI2026} takes a sufficiently large
finite average separately at each stage.  We now give conditions under which
one averaging count works at every stage, together with explicit bounds for
that count.  The main point is to control the error relative to the product
of the interval probabilities.  The product itself will not enter the required
normal-approximation accuracy.  Positive-power growth can also be obtained
nonquantitatively from a Boolean seed by the composition argument in
Section~\ref{subsec:composition}; the result here supplies fixed
parameters and dimension bounds.

For a finite observation $F$ of $N$ independent fair signs, retain the notation
\[
 H_N=\sum_{i=1}^N X_i,\qquad T_F=\E[H_N\mid F],\qquad
 v(F)=\E T_F^2.
\]
When $v(F)>0$, write
\[
 \kappa(F)=\frac{\E T_F^4}{v(F)^2}.
\]
All degree bounds in this section are bounds on the indicators of the
attained cells.  A function of an observation then has the same degree bound.

\subsection{The reporting identity for nonsingleton intervals}

We restate the reporting construction of \cite{OpenAI2026} in the form
needed for quantitative estimates.  Fix $m\ge3$, bounded intervals
$I_1,\ldots,I_m$, a selector $\iota:\R\to\{1,\ldots,m\}$, and a real
number $c$.  Apply the rule to independent copies
$(U,Z_1,\ldots,Z_m)$ of a centered variance-one random variable.  Put
\[
 L_0=U+\sum_{j=1}^m Z_j,\qquad
 \mathsf B=\bigcap_j\{Z_j\in I_j\},\qquad
 \mathsf P=\bigcap_{j\ne\iota(U)}\{Z_j\in I_j\},\qquad
 \mathsf A=\mathsf P\setminus\mathsf B.
\]
The observation $\mathcal W$ reports the tag $\mathsf A$ on $\mathsf A$,
the tag $\mathsf B$ and all leaf values on $\mathsf B$, and the tag
$\mathsf P^c$, the central value, the selected index, and all unselected
leaf values on $\mathsf P^c$.  The identity
\begin{equation}
 \1_{\mathsf A}
 =\sum_{i=1}^m\1_{\{\iota(U)=i\}}
                 \prod_{j\ne i}\1_{I_j}(Z_j)
       -\prod_{j=1}^m\1_{I_j}(Z_j)
 \label{amp:cell-identity}
\end{equation}
shows that every cell indicator is a linear combination of functions of at
most $m$ input coordinates.  For an ordinary report on $\mathsf P^c$, the
reported leaves already certify a failed test, so the omitted leaf is
unrestricted.  This verifies the same assertion for all ordinary reports,
including on finite alphabets with atoms at interval endpoints.

For a centered variance-one law $V$, define
\begin{align}
 q_i^V&=\Pbb(V\in I_i),&
 \mu_i^V&=\E[V\mid V\in I_i],\notag\\
 s_i^V&=\Var(V\mid V\in I_i),&
 Q_V&=\prod_iq_i^V,\quad S_V=\sum_i s_i^V,
 \label{amp:interval-parameters}
\end{align}
whenever all $q_i^V$ are positive.  Thus $s_i^V$ denotes a variance, not
a standard deviation.  With $C_i=\{t:\iota(t)=i\}$, put
\begin{equation}
 J_V=\sum_{i=1}^m\frac1{q_i^V}
       \E\left[\1_{C_i}(V)
       \left((V+c-\mu_i^V)^2+\sum_{j\ne i}s_j^V\right)\right].
 \label{amp:J}
\end{equation}

\begin{proposition}[Exact predictor gain]
\label{amp:reporter}
For the reporting rule above, use the predictor
$b_V=\sum_i\mu_i^V-c$ on $\mathsf A$, the sum of the leaves on
$\mathsf B$, and the sum of the reported values on $\mathsf P^c$.
Its mean squared error is
\begin{equation}
 1-\Delta_V,\qquad
 \Delta_V=Q_V(c^2+S_V-J_V).
 \label{amp:relative-gain}
\end{equation}
Consequently
\begin{equation}
 \Var\bigl(\E[L_0\mid\mathcal W]\bigr)\ge m+\Delta_V.
 \label{amp:variance-gain}
\end{equation}
On every product of finite input alphabets, the observation has finite
range and the coordinate-degree property stated after
\eqref{amp:cell-identity}.
\end{proposition}

\begin{proof}
The error on $\mathsf B$ is the unreported central input.  The error on
each ordinary report is its unreported leaf.  Independence therefore makes
the integrated squared errors on these events equal to their probabilities.
It follows that the total error is
$1-\E[\1_{\mathsf A}(1-(L_0-b_V)^2)]$.

Conditional on $U=t$ and $\mathsf P$, the event has probability
$Q_V/q_{\iota(t)}^V$, and $L_0-b_V$ has mean
$t+c-\mu_{\iota(t)}^V$ and variance
$1+\sum_{j\ne\iota(t)}s_j^V$.  Thus the contribution on
$\mathsf P$ to the last gain is $-Q_VJ_V$.  On $\mathsf B$, the
central input remains unrestricted; the same difference has mean $c$ and
variance $1+S_V$.  Its contribution is $-Q_V(c^2+S_V)$.
Subtracting, since $\1_{\mathsf A}=\1_{\mathsf P}-\1_{\mathsf B}$,
proves \eqref{amp:relative-gain}.  Conditional expectation minimizes
mean squared error, and $\Var(L_0)=m+1$, proving
\eqref{amp:variance-gain}.  The finite-alphabet assertion follows from
\eqref{amp:cell-identity} and the description of the ordinary reports.
\end{proof}

The law-dependent predictor $b_V$ is only a tool for estimating the error.
The observation $\mathcal W$, its cells, and their degree bounds do not
depend on $b_V$.

\subsection{Restricted moments from a Kolmogorov bound}

Write $d_{\mathrm K}$ for Kolmogorov distance and let $G$ denote a standard
normal random variable.

\begin{lemma}[Restricted first and second moments]
\label{amp:restricted}
Suppose $U$ is centered, has variance one, and satisfies $\E U^4\le K$,
where $K\ge3$.  If $d_{\mathrm K}(U,G)\le\delta\le1$ and $E$ is a
union of at most $\ell$ intervals, then
\begin{align}
 |\Pbb(U\in E)-\Pbb(G\in E)|&\le2\ell\delta,\\
 |\E[U\1_E(U)]-\E[G\1_E(G)]|
   &\le(2\ell+3)K^{1/4}\delta^{3/4},\\
 |\E[U^2\1_E(U)]-\E[G^2\1_E(G)]|
   &\le(2\ell+4)\sqrt{K\delta}.
 \label{amp:restricted-bounds}
\end{align}
Unbounded intervals and arbitrary endpoint conventions are allowed.
For one bounded interval $I\subseteq[-B,B]$, the unnormalized moments
also satisfy the sharper estimates
\begin{equation}
 \begin{aligned}
 |q_U-q_G|&\le2\delta,\\
 |\E[U\1_I(U)]-\E[G\1_I(G)]|&\le4B\delta,\\
 |\E[U^2\1_I(U)]-\E[G^2\1_I(G)]|&\le4B^2\delta.
 \end{aligned}
 \label{amp:bounded-moments}
\end{equation}
\end{lemma}

\begin{proof}
Let $D(t)=F_U(t)-\Phi_{\mathrm G}(t)$, where $\Phi_{\mathrm G}$ is
the standard normal distribution function. Both fourth moments are at most $K$.
Markov's inequality, applied to the appropriate right or left tail, gives
\[
 |D(t)|\le\min\{\delta,K/|t|^4\}.
\]
The same bound holds at left limits.  If $\delta>0$, split the following
integrals at $R=(K/\delta)^{1/4}$ to obtain
\begin{equation}
 \int_{\R}|D(t)|\,dt\le\frac83K^{1/4}\delta^{3/4},\qquad
 \int_{\R}2|tD(t)|\,dt\le4\sqrt{K\delta}.
 \label{amp:cdf-integrals}
\end{equation}
At any endpoint, $|tD(t)|\le K^{1/4}\delta^{3/4}$ and
$|t^2D(t)|\le\sqrt{K\delta}$.  Stieltjes integration by parts over
the disjoint intervals now gives the first assertion.  There are at most
$2\ell$ finite endpoint contributions; the integral terms in
\eqref{amp:cdf-integrals} need be counted only once.  The constants
$2\ell+3$ and $2\ell+4$ follow.  Boundary terms at infinity vanish
by the fourth-moment bound.  Left limits handle any atoms at an included
or excluded endpoint.  When $\delta=0$, the laws coincide.

For a bounded interval, the total variations of its indicator, of $t$
times its indicator, and of $t^2$ times its indicator are bounded by
$2$, $4B$, and $4B^2$, respectively.  The same integration-by-parts
argument gives \eqref{amp:bounded-moments}.
\end{proof}

\subsection{Transfer of a fixed Gaussian gain}

\begin{theorem}[Explicit relative transfer]
\label{amp:transfer}
Assume $I_i\subseteq[-B,B]$ with $B\ge1$, and each selector cell $C_i$
is a union of at most $\ell$ intervals.  In the Gaussian model, suppose
\[
 q=\min_i q_i^G>0,\qquad
 \Gamma=c^2+S_G-J_G>0.
\]
For $K\ge3$, set $C=|c|+B$ and
\begin{align}
 A={}&\frac{36mB^2}{q}\notag\\
 &+\frac{2m}{q}\left((2\ell+4)\sqrt K
      +2C(2\ell+3)K^{1/4}+2\ell(C^2+mB^2)\right)\notag\\
 &+\frac{48B(1+C)+76mB^2+4(1+C)^2}{q^2},
 \label{amp:A}\\
 \delta_*={}&\min\left\{1,\frac{q}{4m},
                           \left(\frac{\Gamma}{2A}\right)^2\right\}.
 \label{amp:delta-star}
\end{align}
Every centered variance-one law $U$ with $\E U^4\le K$ and
$d_{\mathrm K}(U,G)\le\delta_*$ satisfies
\begin{equation}
 Q_U\ge Q_G/2,\qquad c^2+S_U-J_U\ge\Gamma/2,\qquad
 \Delta_U\ge Q_G\Gamma/4.
 \label{amp:transfer-conclusion}
\end{equation}
\end{theorem}

\begin{proof}
Write $\delta=d_{\mathrm K}(U,G)$.  Since $\delta\le q/(4m)$,
all interval probabilities for $U$ are positive.  The bounded-interval
estimates in Lemma~\ref{amp:restricted} imply
\begin{align}
 |\mu_i^U-\mu_i^G|&\le\frac{12B\delta}{q},&
 |s_i^U-s_i^G|&\le\frac{36B^2\delta}{q},\\
 \left|\frac1{q_i^U}-\frac1{q_i^G}\right|&\le\frac{4\delta}{q^2},&
 \frac1{q_i^U}&\le\frac2q.
 \label{amp:parameter-errors}
\end{align}
For example, the first conditional moment changes by at most
$(4B\delta+2B\delta)/(q/2)=12B\delta/q$.  The conditional
second moment changes by at most $12B^2\delta/q$; subtracting the
squared means costs at most a further $24B^2\delta/q$.
Moreover,
\[
 \frac{Q_U}{Q_G}
 \ge\prod_i\left(1-\frac{2\delta}{q_i^G}\right)
 \ge1-\frac{2m\delta}{q}\ge\frac12.
\]

For completeness, we bound $J_U-J_G$ with all factors displayed.
Put
\[
 N_i^V(t)=(t+c-\mu_i^V)^2+\sum_{j\ne i}s_j^V.
\]
In \eqref{amp:J}, first replace $N_i^U$ by $N_i^G$ under the law
of $U$, then replace that law by the Gaussian law, and finally replace
$1/q_i^U$ by $1/q_i^G$.  The three absolute errors are at most,
respectively,
\begin{align}
 E_1&=\frac{\delta}{q^2}
                 \left(48B(1+C)+72mB^2\right),\\
 E_2&=\frac{2m}{q}\Bigl((2\ell+4)\sqrt{K\delta}
       +2C(2\ell+3)K^{1/4}\delta^{3/4}\notag\\
    &\hspace{37mm}+2\ell(C^2+mB^2)\delta\Bigr),\\
 E_3&=\frac{4\delta}{q^2}\left((1+C)^2+mB^2\right).
 \label{amp:three-errors}
\end{align}
For $E_1$, use
\[
 |(t+c-\mu_i^U)^2-(t+c-\mu_i^G)^2|
 \le\frac{24B\delta}{q}(|t|+C),
\]
together with \eqref{amp:parameter-errors}.  The cells $C_i$ partition
the line and $\E|U|\le1$, so these terms are summed before taking
expectations.  For $E_2$, expand the fixed quadratic $N_i^G$ and use
\eqref{amp:restricted-bounds}; its three coefficient bounds are
$1$, $2C$, and $C^2+mB^2$.  For $E_3$, use the denominator estimate
and $\E|G|\le1$, again summing over the partition first.

Finally $|S_U-S_G|\le36mB^2\delta/q$.  Since
$\delta\le1$, equations \eqref{amp:A} and \eqref{amp:three-errors}
give
\[
 |S_U-S_G|+|J_U-J_G|\le A\sqrt\delta\le\Gamma/2.
\]
Together with the product bound, this proves
\eqref{amp:transfer-conclusion}.
\end{proof}

The hypotheses are fixed one-dimensional estimates.  In particular,
$\delta_*$ depends on $q$, $\Gamma$, the number of leaves, and the
geometry of the intervals, while the small product $Q_G$ determines
the eventual gain rather than the accuracy requirement.

\subsection{One averaging count for every stage}

\begin{theorem}[Fixed-count amplification]
\label{amp:fixed-count}
Assume the hypotheses of Theorem~\ref{amp:transfer}, use $K=6$, and
choose an integer
\[
 M\ge\max\{2,\lceil6/\delta_*^2\rceil\}.
\]
Let $0<\eta\le Q_G\Gamma/4$, and set
\begin{equation}
 \begin{aligned}
 d_*&=mM,\qquad n_*=(m+1)M,\qquad v_*=M(m+\eta),\\
 \alpha&=\frac{\log v_*}{2\log d_*}
       =\frac12+\frac{\log(1+\eta/m)}{2\log(mM)}.
 \end{aligned}
 \label{amp:scales}
\end{equation}
Writing $L(f)=\sum_i\widehat f(\{i\})$, there are finite observations
$F_k$ and Boolean readouts $f_k$ such that
\begin{align}
 N_k&=n_*^k,& \degc(F_k)&\le d_*^k,\notag\\
 v(F_k)&\ge v_*^k,& \kappa(F_k)&\le6,\notag\\
 \deg(f_k)&\le d_*^k,& L(f_k)&\ge6^{-1/2}d_*^{k\alpha}.
 \label{amp:iterates}
\end{align}
In particular, the Boolean degrees tend to infinity and
$L(f_k)\ge6^{-1/2}\deg(f_k)^\alpha$.
\end{theorem}

\begin{proof}
For a finite observation with score $T$, positive variance $v$, and
$\E T^4/v^2\le6$, take $M$ independent copies and form
\[
 Z=\frac{T_1+\cdots+T_M}{\sqrt{Mv}}.
\]
Its score as an observation of the underlying signs is exactly
$\sqrt{Mv}\,Z$, and its cell degree is at most $M\degc(F)$.
Writing $\kappa=\E T^4/v^2$, independence gives
\begin{equation}
 \E Z^4=3+\frac{\kappa-3}{M}.
 \label{amp:average-fourth}
\end{equation}
The Berry--Esseen inequality with an absolute constant below one
\cite{Shevtsova2011}, and
$\E|T/\sqrt v|^3\le\sqrt\kappa$, give
\begin{equation}
 d_{\mathrm K}(Z,G)\le\sqrt{6/M}\le\delta_*.
 \label{amp:BE}
\end{equation}
Thus Theorem~\ref{amp:transfer} applies to $m+1$ independent copies
of $Z$.  If $F'$ is their report, the conditional-expectation identity
for the total underlying sign sum gives
\[
 v(F')\ge Mv(m+\eta),\qquad \degc(F')\le mM\degc(F).
\]

It remains to check the fourth moment uniformly.  Let $r=m+1$,
$S_0=Z_1+\cdots+Z_r$, and
$V=\Var(\E[S_0\mid F'])\ge m+\eta$.  Conditional Jensen's
inequality and independence yield
\begin{align}
 \kappa(F')\le\frac{\E S_0^4}{V^2}
 &\le\frac{3r^2+r(\kappa-3)/M}{(m+\eta)^2}\notag\\
 &\le\frac{3(m+1)^2+3(m+1)/2}{m^2}\le6.
 \label{amp:moment-invariant}
\end{align}
The last inequality follows from
\[
 6m^2-3(m+1)^2-\tfrac32(m+1)
 =\tfrac32(2m+1)(m-3)\ge0.
\]
Start with one revealed sign, whose variance and cell degree are both
one and whose standardized fourth moment is one.  Induction proves
the observation assertions in \eqref{amp:iterates} with the same $M$
at every step.

Take $f_k=\operatorname{sign}(T_{F_k})$, with sign $+1$ at zero.
It is a function of the observation, so has the stated degree bound,
and conditioning gives $L(f_k)=\E|T_{F_k}|$.
For any centered variance-one $Y$, H\"older's inequality gives
\[
 1=\E Y^2\le(\E|Y|)^{2/3}(\E Y^4)^{1/3}.
\]
Thus $\E|T_{F_k}|\ge\sqrt{v(F_k)/6}$, which proves the lower
bound.  Finally every Boolean function satisfies $L(f)\le\deg(f)$:
each absolute singleton coefficient is at most its influence, and the
sum of the influences is at most the degree by Parseval.  Hence the
degrees here tend to infinity.
\end{proof}

\begin{corollary}[Degree budgets and ambient dimension]
\label{amp:dimension}
With the notation of Theorem~\ref{amp:fixed-count}, put
\[
 \theta=\frac{\log n_*}{\log d_*},\qquad
 \beta=\frac{2\log n_*}{\log v_*}.
\]
For every real degree budget $D\ge d_*$, one of the readouts has
\begin{equation}
 \deg(f)\le D,\qquad n\le D^\theta,\qquad
 L(f)\ge6^{-1/2}d_*^{-\alpha}D^\alpha.
 \label{amp:budget}
\end{equation}
For each individual iterate, its actual degree $d=\deg(f_k)$ also gives
\begin{equation}
 N_k\le(\sqrt6\,d)^\beta.
 \label{amp:actual-degree}
\end{equation}
\end{corollary}

\begin{proof}
Choose $k=\lfloor\log D/\log d_*\rfloor$.  Then
$D/d_*<d_*^k\le D$, and \eqref{amp:budget} follows from
\eqref{amp:iterates}.  For \eqref{amp:actual-degree}, the already
proved inequality
$d\ge L(f_k)\ge6^{-1/2}v_*^{k/2}$ implies
$(\sqrt6\,d)^\beta\ge n_*^k=N_k$.
\end{proof}

If all interval endpoints and selector cut points are rational, these
constructions are terminating exact algorithms.  Every unnormalized score
and its variance are rational because the underlying cube is finite.  A
normalized score is a quadratic algebraic number, so comparisons with the
fixed rational cut points are decidable, including equality.  Conditional
scores can be obtained by finite enumeration of cells.  The dimension and
degree bounds quantify mathematical size; no efficient running-time claim
is made.
